\chapter{Worked example: the Cryptobox game-hopping proof}

NaCl's \texttt{crypto\_box} combines a non-interactive key exchange (NIKE) with a
nonce-based symmetric encryption scheme to obtain public-key authenticated
encryption (PKAE). The module set follows Dupressoir et al., \emph{Bringing State
Separating Proofs to EasyCrypt} (CSF 2022), and its Rocq SSProve counterpart. The
security theorem is a chain of game hops whose individual steps are either
proved or exposed as named hypotheses; this chapter records which is which.

\begin{definition}[Schemes and games]
  \label{def:cryptobox-games}
  \uses{def:spcomp, def:advantage}
  \lean{CatCrypt.Examples.Cryptobox.NBSES, CatCrypt.Examples.Cryptobox.NIKEScheme, CatCrypt.Examples.Cryptobox.PKEY, CatCrypt.Examples.Cryptobox.KEY, CatCrypt.Examples.Cryptobox.SAE, CatCrypt.Examples.Cryptobox.GAE, CatCrypt.Examples.Cryptobox.GNIKE, CatCrypt.Examples.Cryptobox.PKAEGame.GPKAE, CatCrypt.Examples.Cryptobox.PKAE}
  \leanok
  The nonce-based symmetric scheme, the NIKE, and the stateful games for public-key
  management (PKEY), shared keys (KEY), single- and multi-instance authenticated
  encryption (SAE, AE), NIKE security, and PKAE security, each with a real and an
  ideal variant.
\end{definition}

\begin{theorem}[PKEY switching steps]
  \label{thm:cryptobox-pkey}
  \uses{def:cryptobox-games, thm:advantageA-triangle}
  \lean{CatCrypt.Examples.Cryptobox.PKAE_eq_PKEY_nocheck, CatCrypt.Examples.Cryptobox.pkae_game_hopping}
  \leanok
  The PKAE game and its collision-checking variant are the PKEY game composed with
  an explicit reduction adversary (a definitional equality), so the PKAE
  advantage is at most two PKEY advantages plus the advantage between the real
  and ideal collision-checking games.
\end{theorem}

\begin{theorem}[Collision bound for PKEY]
  \label{thm:cryptobox-collision}
  \uses{def:cryptobox-games}
  \lean{CatCrypt.Examples.Cryptobox.pkeyAdvantage_le_collision, CatCrypt.Examples.Cryptobox.pkey_switching, CatCrypt.Examples.Cryptobox.collision_prob_multi_query_bound}
  \leanok
  If the two PKEY games agree on every run that does not set the collision flag
  and the real game is lossless, the PKEY advantage is at most the probability of
  the collision flag, and hence at most any bound on that probability. The sum of the per-query collision probabilities over
  $q_{\mathrm{gen}}$ generation queries is computed in closed form.
\end{theorem}

\begin{theorem}[AE hybrid over sessions]
  \label{thm:cryptobox-ae}
  \uses{def:cryptobox-games, thm:advantageA-triangle}
  \lean{CatCrypt.Examples.Cryptobox.aeHybrid_empty, CatCrypt.Examples.Cryptobox.aeHybrid_univ, CatCrypt.Examples.Cryptobox.aeStep_triangle, CatCrypt.Examples.Cryptobox.aeStep_bounded_by_sae}
  \leanok
  The AE switching step is a hybrid over session identifiers with the stated
  boundary games. Under the hypothesis that each hybrid step is the SAE game run
  against an explicit reduction adversary, its advantage is at most the sum of
  the per-session SAE bounds.
\end{theorem}

\begin{theorem}[Cryptobox security]
  \label{thm:cryptobox}
  \uses{thm:cryptobox-pkey, thm:cryptobox-collision, thm:cryptobox-ae}
  \lean{CatCrypt.Examples.Cryptobox.cryptobox_security_chain, CatCrypt.Examples.Cryptobox.cryptobox_security_full, CatCrypt.Examples.Cryptobox.cryptobox_security_with_nike_reduction}
  \leanok
  Given bounds $\varepsilon_{\mathrm{pkey}}$ on both PKEY reduction advantages,
  $\varepsilon_{\mathrm{nike}}$ on the NIKE switching step and
  $\varepsilon_{\mathrm{ae}}$ on the AE switching step, the PKAE advantage is at
  most $2\varepsilon_{\mathrm{pkey}} + \varepsilon_{\mathrm{nike}} +
  \varepsilon_{\mathrm{ae}}$. The last variant replaces the NIKE step by the
  advantage of an explicit NIKE reduction plus a caching gap that is a
  hypothesis.
\end{theorem}
